% Chapter 4
\chapter{Data}
\label{Chapter4} % For referencing the chapter elsewhere, use \ref{Chapter3} 

This chapter describes the datasets used to construct the financial-environmental attribution framework. The analysis combines three main data components. First, from-whom-to-whom financial accounts are used to measure direct and indirect financial exposures across institutional sectors. Second, environmental indicators from EXIOBASE are used to quantify the environmental pressures generated by real-economy agents. Third, balance-sheet data are used to construct the attribution base against which financial exposures are normalized. The final section explains the exchange-rate conversion used to express all monetary variables in a common currency.

\section{From-whom-to-whom matrices}

\textit{FWTW data in the United States. } The Federal Reserve provides FWTW data for each financial instrument, available at \url{https://www.federalreserve.gov/releases/efa/fwtw.htm}
. These data are constructed to align with the issuer, holder, and instrument totals reported in the Financial Accounts. The series is available on a quarterly basis from 1945 to the most recent quarter. 

Table \ref{tab:fwtw_classification} summarizes the sources used to construct the FWTW data for each financial instrument in the United States. As of the fourth quarter of 2022, approximately two-thirds of total asset values are fully specified by \textit{“known relationships”} and \textit{“exclusion restrictions”}, while the remainder relies partially or entirely on proportionality assumptions.

\begin{table}[ht]
\centering
\caption{FWTW data by financial instrument in the US}
\footnotesize
\setlength{\tabcolsep}{4pt}
\renewcommand{\arraystretch}{0.9}

\begin{tabularx}{\textwidth}{>{\raggedright\arraybackslash}p{0.28\textwidth} 
                                  >{\raggedright\arraybackslash}X 
                                  >{\raggedright\arraybackslash}X}
\toprule
\multicolumn{3}{c}{\textbf{Known relationships and exclusion restrictions fully specify FWTW data (\$171.6 trillion as of 2022:Q4)}} \\
\midrule

\multicolumn{3}{l}{
\begin{minipage}{\textwidth}
\vspace{1pt}
\begin{itemize}[leftmargin=*, itemsep=0pt, topsep=0pt]
\item U.S. Official Reserve Assets and SDR Allocations
\item SDR Certificates and Treasury Currency
\item U.S. Deposits in Foreign Countries
\item Net Interbank Transactions
\item Time and Savings Deposits
\item Money Market Funds Shares
\item Treasury Securities
\item Depository Institution Loans N.E.C.
\item Other Loans and Advances
\item Farm Mortgages
\item Consumer Credit
\item Mutual Fund Shares
\item Life Insurance Reserves
\item Pension Entitlements
\item Proprietors' Equity in Noncorporate Business
\item Direct Investment
\item Identified Miscellaneous Financial Claims - I
\item Identified Miscellaneous Financial Claims - II
\end{itemize}
\vspace{1pt}
\end{minipage}
} \\

\midrule

 & \textbf{Single-sector dominated} & \textbf{Multi-sector balanced} \\
\midrule

\textbf{Known relationships, restrictions, and proportionality assumptions} 
& (\$19.4T as of 2022:Q4)
\begin{itemize}[leftmargin=*, itemsep=0pt, topsep=0pt]
\item Municipal Securities
\item Home Mortgages
\item Multifamily Residential Mortgages
\end{itemize}
& (\$33T as of 2022:Q4)
\begin{itemize}[leftmargin=*, itemsep=0pt, topsep=0pt]
\item Checkable Deposits and Currency
\item Federal Funds and Repos
\item Open Market Paper
\item Corporate and Foreign Bonds
\item Taxes Payable
\end{itemize} \\

\midrule

\textbf{Proportionality assumptions only} 
& (\$15.3T as of 2022:Q4)
\begin{itemize}[leftmargin=*, itemsep=0pt, topsep=0pt]
\item Agency- and GSE-Backed Securities
\item Commercial Mortgages
\end{itemize}
& (\$19.4T as of 2022:Q4)
\begin{itemize}[leftmargin=*, itemsep=0pt, topsep=0pt]
\item Trade Credit
\item Unidentified Miscellaneous Financial Claims
\end{itemize} \\

\bottomrule
\end{tabularx}

\vspace{4pt}
\begin{minipage}{\textwidth}
\footnotesize
\textit{Source: The Federal Reserve}
\end{minipage}

\label{tab:fwtw_classification}
\end{table}

Raw FWTW data in the United States are provided with information on instrument names (and codes), holder names (and codes), issuer names (and codes), as well as the corresponding date and level. This thesis uses fourth-quarter observations spanning from 2000 to 2024. 

Following the classification outlined in Section \ref{Section 3.1.2}, holders and issuers are aggregated into four real-economy sectors and five financial institution categories. To account for discrepancies, an additional row and column are introduced, resulting in a $10 \times 10$ matrix for each financial instrument. These instruments are then consolidated into ten major asset classes, and the matrix construction procedure is repeated annually over the sample period.

% ---------------
\textit{FWTW data in the euro area.} Compared to the United States, the processing and construction of FWTW matrices in the Euro area is more complex. Financial accounts and FWTW-related data for EU member states and the Euro area are published by the ECB within a comprehensive dataset known as the \textit{Quarterly Sector Accounts}, which contains 369{,}626 time series. To navigate this dataset, it is necessary to consult the accompanying catalogue, which provides a complete list of series and their associated metadata and is available at \href{https://data.ecb.europa.eu/data/data-categories/2045085/data-information?layerType=KL&showDatasetModal=false}{this link}.

For each financial instrument, the catalogue is filtered to retain only quarterly euro area series (\texttt{ref\_area = I9}), reported on an end-of-period basis (\texttt{time\_per\_collect = E}), and at the total level of breakdown (\texttt{cust\_breakdown = \_T}). For loans and deposits, which are further disaggregated by maturity in the ECB database, only total-maturity series (\texttt{maturity = T}) are retained to avoid duplication. The dataset is further restricted to series for which both the reporting sector (\texttt{ref\_sector}) and the counterpart sector (\texttt{counterpart\_sector}) belong to the set of institutional sectors defined in the previous classification.

Once the relevant series have been identified, the associated series codes are used to query the ECB Data API and retrieve the corresponding observations, which are then combined by instrument and exported as separate CSV files. Each series includes a field indicating whether the position is recorded as an asset or a liability from the perspective of the reporting sector; this information is used to identify the holder and issuer of the corresponding instrument.

The ECB provides complete FWTW matrices for a subset of instruments, including deposits, debt securities, loans, listed shares, money market fund shares, and
mutual fund shares. However, the time coverage differs across instruments: for some instruments, complete FWTW information is available only from 2013 onward.
For the remaining instruments and periods, the ECB reports only aggregate row and column totals through the financial accounts. To address this limitation, I
define ``restriction exclusions'' and identify available ``known relationships.'' I then apply the IPFP algorithm to estimate and populate the missing interior
entries of the matrix. Table~\ref{tab:fwtw_eu} summarizes which instruments are directly sourced from the ECB and which are constructed using this methodology.
The structural restrictions imposed for each financial instrument are presented in Appendix~\ref{Appendix A: Structural Restrictions}.

As discussed in Section~\ref{sec:FWTW_introduction}, total issuance may not equal total holdings for some financial instruments. In such cases, the algorithm is set to match the row sums, i.e., the total holdings of each sector for the corresponding instrument. The main reason for prioritizing row balancing over column balancing is to ensure comparability with US FWTW data, where a similar approach is adopted, resulting in a zero discrepancy (DIS) column. To reconcile remaining inconsistencies, discrepancy (DIS) rows and columns are introduced so that the matrices are consistent with aggregate totals reported in the financial accounts.
\begin{table}[ht!]
\centering
\caption{FWTW data by financial instrument in the Euro area}
\label{tab:fwtw_eu}
\small
% Note: \textwidth makes it fit exactly to the page margins
\begin{tabularx}{\textwidth}{@{} 
    >{\raggedright\arraybackslash}p{3cm} 
    >{\raggedright\arraybackslash}X 
    l 
    >{\centering\arraybackslash}p{2.5cm} 
    >{\centering\arraybackslash}p{2.5cm} 
    @{}}
\toprule
\textbf{Group} & \textbf{Instruments} & \textbf{Code} & \textbf{FWTW by ECB} & \textbf{FWTW by author} \\
\midrule
% ---------------------------------------------------------
\textbf{Monetary Gold and SDRs} & Monetary gold and SDRs & F1 & & \checkmark \\
\midrule
% ---------------------------------------------------------
\textbf{Currency and Deposits} & Currency & F2M &  & \checkmark \\
& Deposits & F21 & \checkmark &  \\
\midrule
% ---------------------------------------------------------
\textbf{Debt securities} & Debt securities (*) & F3 & \checkmark  & \checkmark \\
 & & & (2013-2024) & (2000-2012) \\
\midrule
% ---------------------------------------------------------
\textbf{Loans} & Loans & F4 & \checkmark & \\
\midrule
% ---------------------------------------------------------
\textbf{Equity} & Listed shares & F511 & \checkmark & \checkmark\\
 & & & (2013-2024) & (2000-2012) \\
& Unlisted shares & F512 & & \checkmark \\
& Other equity & F519 & & \checkmark \\
\midrule
% ---------------------------------------------------------
\textbf{Money Market Funds Shares} & Money market fund shares (**) & F521 & \checkmark & \checkmark \\
 & & & (2013-2024) & (2000-2012) \\
\midrule
% ---------------------------------------------------------
\textbf{Mutual Fund Shares} & Non-MMF investment fund shares (**) & F522 & \checkmark & \checkmark \\
 & & & (2013-2024) & (2000-2012) \\
\midrule
% ---------------------------------------------------------
\textbf{Insurance} & Life insurance & F62 & & \checkmark \\
& Non-life insurance technical reserves and provision for calls under standardised guarantees & F6O & & \checkmark \\
\midrule
% ---------------------------------------------------------
\textbf{Pension} & Pension schemes & F6M & & \checkmark \\
\midrule
% ---------------------------------------------------------
\textbf{Other} & Financial derivatives & F7 & & \checkmark \\
& Other accounts receivable & F8 & & \checkmark \\
\bottomrule
\end{tabularx} 

\vspace{4pt}
\begin{minipage}{\textwidth}
\footnotesize
\textit{Source: Author's compilation.}\\

(*) The ECB provides the full matrix for debt securities from 2013 onward. Before 2013, only row and column totals are available from the financial accounts. The ECB also reports holdings of debt securities issued by and held within non-financial corporations, which is used as a known relationship in the estimation.\\

(**) All sectors hold money market fund (MMF) and mutual fund (non-MMF) shares, but issuance is limited to RoW and MFI (for MMF shares) and NM\_MFIF (for non-MMF shares). From 2013 onwards, ECB reports, for each sector, holdings of shares issued by MFI (and NM\_MFIF), but provides only total issuance for RoW. As there are only two issuers, the RoW column is derived residually as total holdings minus the MFI (or NM\_MFIF) column. 
\end{minipage}
\end{table}

\textit{Discrepancies.} While it is important to verify that each component sums to the aggregate macroeconomic totals, discrepancies do not have economic substance as either real-economy agents or financial institutions. Accordingly, after aggregating the FWTW matrices across all asset classes (see Section~\ref{sec:effective_financing}), I exclude the discrepancy row and column from subsequent steps when computing effective exposures. Although discrepancies can account for a substantial share of issuance or holdings for certain financial instruments (e.g., monetary gold and SDRs), their aggregate importance remains limited. In particular, total discrepancies remains below 3 percent in the United States and below 1 percent in the Euro area of total issuance across all financial assets for any years in the 2000-2022 period. Detailed figures are presented in Appendix \ref{Appendix B:Total_dis}.

\section{Environmental indicators}

To quantify the environmental pressures generated by real-economy agents, I utilize EXIOBASE v3, an \textit{environmentally extended multi-regional input-output} (EE-MRIO) database that captures the global web of economic relationships and their corresponding environmental impacts \citep{stadler_exiobase_2018}. Input-output modelling is a well-established framework that traces monetary flows between industries and end-users to represent global production and consumption structures. By integrating environmental satellite accounts, this framework enables the quantification of pressures such as greenhouse gas emissions, land and water use, and material extraction, tracing them along international supply chains \citep{hauschild_life_2018, leclerc_building_2020}. EXIOBASE v3 is one of the most comprehensive EE-MRIO databases available, covering 44 individual countries and five RoW regions, disaggregated into 163 industries, with an annual time series spanning 1995 to 2022.

Although the attribution framework detailed in Chapter \ref{Chapter3} can accommodate any environmental indicators available in EXIOBASE v3, this thesis focuses on four key dimensions of environmental impact. For each dimension, I extract direct production pressures ($F_t$) and direct final-use pressures ($F_{Y,t}$) from the respective EXIOBASE impact accounts:

\begin{itemize}
    \item \textbf{Climate Change:} Greenhouse gas emissions, captured via the \texttt{GHG emissions GWP100 IPCC2010} account. This aggregates carbon dioxide (CO$_2$), methane (CH$_4$), and nitrous oxide (N$_2$O) into kilograms of carbon dioxide equivalents (kg CO$_2$e) using the 100-year Global Warming Potentials established by the IPCC.
    \item \textbf{Land Use:} Total land occupation, measured in square kilometers (km$^2$).
    \item \textbf{Water Consumption:} Blue water consumption, measured in millions of cubic meters (Mm$^3$).
    \item \textbf{Resource Extraction:} Domestic extraction of four primary material categories, including fossil fuels, iron ore, non-ferrous metal ores, and non-metallic minerals, measured in kilotonnes (kt).
\end{itemize}

Choosing EXIOBASE v3 over other global EE-MRIO databases, such as Eora, the OECD Inter-Country Input-Output Tables (ICIO), and the World Input-Output Database (WIOD), involves several methodological trade-offs.

\textit{Advantages.} EXIOBASE v3 is particularly suitable for this thesis because of the breadth, comparability, and recent coverage of its environmental satellite accounts. While WIOD \citep{timmer_illustrated_2015} and the OECD ICIO \citep{oecd_inter-country_2026} provide reliable economic input-output data, they were primarily designed for macroeconomic and trade analysis. Their environmental extensions are therefore more limited, focusing mainly on energy use, CO$_2$, or GHG emissions. In addition, the most comprehensive WIOD environmental satellite accounts end in 2016, which limits their usefulness for analysing recent macro-financial exposures.

By contrast, EXIOBASE v3 provides annual data up to 2022 and covers a wide range of environmental pressures, including greenhouse gas emissions, land use, water use, energy use, and material extraction \citep{stadler_exiobase_2018}. This makes it possible to conduct a multi-dimensional assessment of the environmental footprint of financial systems, rather than focusing only on climate-related indicators.

EXIOBASE also offers an advantage in terms of sectoral and environmental harmonisation. Eora, for example, provides broader geographic coverage and preserve national input-output tables in their original classifications, but this heterogeneous structure can make cross-country comparisons more difficult and may require additional harmonisation \citep{lenzen_building_2013, casella_unctad_2019}. More generally, differences in classification, aggregation, and database construction choices can lead to non-negligible deviations in footprint estimates across MRIO databases \citep{giljum_impacts_2019}. In this context, the relatively detailed and harmonised structure of EXIOBASE makes it well suited for comparing environmental pressures across the United States and the Euro area.

\textit{Limitations.} Despite its advantages, EXIOBASE presents several limitations relevant to this thesis. First, data for the most recent years are not derived entirely from primary source tables. Although recent EXIOBASE releases extend the time series to 2022, the final years rely on macroeconomic indicators, trade data, and available environmental proxies rather than complete supply-use tables for all countries and sectors \citep{stadler_exiobase_2018}. Consequently, short-term fluctuations in the latest years of the sample, particularly between 2020 and 2022, should be interpreted with caution.

Second, EXIOBASE trades geographic resolution for sectoral and environmental detail. The database explicitly covers 44 countries alongside five aggregate RoW regions. While this structure is well-suited for analysing major economies like the United States and the Euro area, it restricts the ability to capture nuances within developing and emerging economies. Grouping countries with distinct production structures, resource endowments, and environmental intensities into broad RoW regions can obscure substantial local heterogeneity. This is particularly problematic for spatially specific environmental pressures, such as land use, water stress, and biodiversity impacts. Indeed, studies increasing the geographic resolution of EXIOBASE demonstrate that disaggregating RoW regions can materially alter footprint estimates \citep{bjelle_adding_2020, cabernard_highly_2021}.

While these constraints do not undermine EXIOBASE's suitability for this analysis, they must inform how the results are interpreted. Therefore, the estimates presented here should be viewed as harmonised, macro-level footprints of financial systems rather than precise, country-level measures of individual foreign exposures.
\section{Attribution Base}

The attribution base, denoted by $TC_{k,t}$, represents the total claims on real-economy agents $k \in \{\text{NFC}, \text{GG}, \text{RoW}\}$ and the total assets for household sector, where $k = \text{HH}$ (see Section \ref{sec:3.3} for a detailed explanation). For domestic sectors, total claims can be extracted from financial accounts, corresponding to the column sums of the aggregated FWTW matrices. To determine total assets for HH, the net worth of household sector, that is available in financial accounts, is added to total claims on HH. 

For the RoW, the OECD Financial Accounts and Balance Sheets database\footnote{Source: \url{https://www.oecd.org/en/data/datasets/financial-accounts-and-balance-sheets.html}} provides sector-level balance sheets for OECD and several non-OECD countries, serving as a reliable source for indicative figures of the RoW's total claims. Specifically, the total claims of the RoW sector aggregate the liabilities of foreign non-financial corporations, general governments, and financial institutions. The foreign household sector is explicitly excluded from this aggregation, under the assumption that domestic financial institutions do not typically extend direct loans or credit to foreign households. Although this dataset covers 42 countries, only those also covered by the EXIOBASE v3 database are used to construct the RoW attribution base to ensure a consistent and fair attribution factor. By definition, the home country or region (e.g., the US or the Eurozone) is excluded from the calculation of the RoW's total claims. Finally, Appendix \ref{Appendix C:Countries} provides a comprehensive list of countries covered by the OECD and EXIOBASE v3 databases, their regional grouping, the member states of the Eurozone, and the specific countries that constitute the "Rest of World" for the US and Euro area. 

A major advantage of the OECD database is that it is harmonized under the 2008 SNA framework. However, it excludes China, a critical omission given that China holds the second-highest national debt (\cite{world_population_review_countries_2026}) and is the leading global GHG emitter (\cite{european_commission_joint_research_centre_ghg_2025}). To address this gap and align with EXIOBASE v3, which explicitly captures China’s environmental pressures, I incorporate China's financial accounts into the analysis. Specifically, for the broadest RoW scenario, I integrate China's data into the attribution base of the RoW sector for both the US and the Euro area.

In 2020, the People’s Bank of China (PBC) published data on the levels of financial assets and liabilities, i.e., sectoral balance sheets, on its official website for the first time.\footnote{Available at \url{https://www.pbc.gov.cn/en/3688247/3688975/index.html}.} The published series were backfilled to 2017. To bridge the data gap in analyses of China’s macroeconomic operations, academic studies have also used China’s sovereign balance sheets compiled by \textcite{li_yang_chinas_2020} at the Centre for National Balance Sheet, which cover the period 2000-2019.
I compare the overlapping years, 2017-2019, between the two sources (see Table \ref{tab:china_comparison}). The comparison shows that the combined liabilities of NFCs, general government, and financial institutions reported by \textcite{li_yang_chinas_2020} are approximately 35-36\% higher than the corresponding official figures published by the People’s Bank of China in each of the three overlapping years. To ensure consistency with the official PBC series, I therefore adjust the total liabilities of NFCs, general government, and financial institutions from the Centre for National Balance Sheet for the period 2000-2016 downward by the average difference 36.26\%, under the assumption that the structural reporting variance between the academic and official accounts remained relatively stable over the historical period.

\begin{table}[H]
\centering
\renewcommand{\arraystretch}{1.3} 
\caption{China's financial accounts by different sources in \\overlapping years}
\label{tab:china_comparison}

\makebox[\textwidth][r]{\footnotesize Unit: 100 Million Yuan}\vspace{2pt}

\resizebox{\textwidth}{!}{
\begin{tabular}{l r r r r r r r r r}
\toprule
 & \multicolumn{4}{c}{\textbf{The People's Bank of China}} & \multicolumn{4}{c}{\textbf{The Centre for National Balance Sheet}} & \\
\cmidrule(lr){2-5} \cmidrule(lr){6-9}
 & \multicolumn{1}{c}{NFC} & \multicolumn{1}{c}{GG} & \multicolumn{1}{c}{FI} & \multicolumn{1}{c}{\cellcolor{blue!15}NFC+GG+FI} & \multicolumn{1}{c}{NFC} & \multicolumn{1}{c}{GG} & \multicolumn{1}{c}{FI} & \multicolumn{1}{c}{\cellcolor{blue!15}NFC+GG+FI} & \multicolumn{1}{c}{\% Change} \\
 & \multicolumn{1}{c}{\textit{(1)}} & \multicolumn{1}{c}{\textit{(2)}} & \multicolumn{1}{c}{\textit{(3)}} & \multicolumn{1}{c}{\cellcolor{blue!15}\textit{(4) = (1)+(2)+(3)}} & \multicolumn{1}{c}{\textit{(5)}} & \multicolumn{1}{c}{\textit{(6)}} & \multicolumn{1}{c}{\textit{(7)}} & \multicolumn{1}{c}{\cellcolor{blue!15}\textit{(8) = (5)+(6)+(7)}} & \multicolumn{1}{c}{\textit{(9)=(8)/(4) - 1}} \\
\midrule
2017 & 1,875,034 & 416,104 & 3,733,134 & \cellcolor{blue!15}6,024,272 & 3,827,714 & 299,051 & 4,027,658 & \cellcolor{blue!15}8,154,423 & 35.36\% \\
2018 & 1,833,366 & 494,401 & 3,843,805 & \cellcolor{blue!15}6,171,572 & 3,947,967 & 332,666 & 4,160,287 & \cellcolor{blue!15}8,440,920 & 36.77\% \\
2019 & 2,073,460 & 561,773 & 4,081,652 & \cellcolor{blue!15}6,716,885 & 4,392,160 & 379,577 & 4,406,288 & \cellcolor{blue!15}9,178,025 & 36.64\% \\
\midrule
 & & & & & & & & \cellcolor{blue!15}\textbf{\textit{Average}} & 36.26\% \\
\bottomrule
\end{tabular}
}
\vspace{1pt}\\
\begin{minipage}{\textwidth}
\footnotesize\textit{Source: The People's Bank of China and The Centre for National Balance Sheet}
\end{minipage}
\end{table}

\section{Foreign Exchange}

To ensure comparability across datasets, all monetary variables are converted into a common reporting currency, the US dollar. Specifically, end-of-year spot exchange rates are employed to maintain temporal consistency. The Chinese Yuan Renminbi is converted into U.S. dollars using the Chinese Yuan Renminbi to U.S. Dollar spot exchange rate (DEXCHUS)\footnote{Source: \url{https://fred.stlouisfed.org/series/DEXCHUS}}, while values originally denominated in euros are converted using the U.S. Dollar to Euro spot exchange rate (DEXUSEU)\footnote{Source: \url{https://fred.stlouisfed.org/series/DEXUSEU}}. Both series are obtained from the \textit{Federal Reserve Economic Data (FRED)} database maintained by the Federal Reserve.